\section{Related Work}
\label{sec:related}

\subsection{Image-to-Video Diffusion Models}

I2V diffusion models animate a static image by generating
temporally coherent video frames conditioned on the input image
and a text prompt. Open-source backbones such as Wan2.1-I2V
and LTX-Video provide controllable synthesis with manageable GPU
footprints. StyleMaster \citep{ye2025stylemaster} improves
stylized video quality through motion adapters and patch-level
style filtering. EchoStyle \citep{li2026echostyle} uses
sliding-window strategies for long-video temporal coherence.
However, these methods are trained on natural imagery and do not
account for art-specific aesthetic principles such as ink density
or negative space.

\subsection{Ink-Wash Painting Stylization}

Computational ink-wash painting has focused primarily on
\emph{single-image} generation. BitGAN \citep{he2023bitgan}
generates ink-wash images via GANs with brush-texture loss.
ChipDiff \citep{liu2026chipdiff} applies loss-gradient guidance on
a diffusion model to capture ink aesthetics. InkSpirit
\citep{zhang2025inkspirit} uses ControlNet conditioning for
ink-style image transfer. These methods demonstrate that
art-specific losses can guide diffusion models toward ink
aesthetics, but they are limited to static images and do not
address temporal coherence or the five-level ink theory as
explicit controllable parameters.

\subsection{Controllable Diffusion with Art-Specific Conditioning}

ControlNet \citep{zhang2023controlnet} enables structural
conditioning (Canny edges, depth maps, poses) in diffusion models.
IP-Adapter and LoRA \citep{hu2022lora} provide style conditioning
through reference images. However, existing control mechanisms do
not expose art-specific axes such as ink-density levels or
negative-space regions as controllable parameters. Our work fills
this gap by introducing dual conditioning with explicit
art-theoretic foundations and a blank-space loss that has no
counterpart in general-purpose controllable diffusion.

\subsection{Content Deformation Fields for Video Consistency}

CoDeF \citep{ouyang2024codef} introduces content deformation fields
for temporally consistent video-to-video translation. We adapt this
concept for cross-region consistency in multi-region ink-wash
animations, ensuring that motion in one region (e.g., drifting
mist) does not introduce artifacts at boundaries with static
regions (e.g., mountains).
